\documentclass[11pt]{article}

\usepackage[a4paper,margin=2.4cm]{geometry}
\usepackage{fontspec}
\setmainfont{Helvetica Neue}
\usepackage{amsmath}
\usepackage{amssymb}
\usepackage{titlesec}
\usepackage{enumitem}
\usepackage{xcolor}
\usepackage{listings}
\usepackage{booktabs}
\usepackage{hyperref}
\usepackage{fancyhdr}
\usepackage{tcolorbox}
\tcbuselibrary{skins,breakable}

\definecolor{codebg}{RGB}{244,244,247}
\definecolor{codeframe}{RGB}{210,210,216}
\definecolor{accent}{RGB}{31,74,121}
\definecolor{notebg}{RGB}{237,242,247}

\lstset{
  basicstyle=\ttfamily\small,
  breaklines=true,
  backgroundcolor=\color{codebg},
  frame=single,
  rulecolor=\color{codeframe},
  framesep=6pt,
  xleftmargin=2pt,
  xrightmargin=2pt,
  columns=fullflexible,
  keepspaces=true,
}

\newtcolorbox{note}{colback=notebg,colframe=codeframe,boxrule=0.5pt,arc=2pt,left=8pt,right=8pt,top=6pt,bottom=6pt,breakable}

\titleformat{\section}{\Large\bfseries\color{accent}}{\thesection}{0.6em}{}
\titleformat{\subsection}{\large\bfseries\color{accent}}{\thesubsection}{0.6em}{}

\pagestyle{fancy}
\fancyhf{}
\fancyhead[L]{\small\color{gray} Axion Strings AMR --- Physics Reference}
\fancyhead[R]{\small\color{gray} \thepage}
\renewcommand{\headrulewidth}{0.3pt}

\hypersetup{
  colorlinks=true,
  linkcolor=accent,
  urlcolor=accent,
  citecolor=accent,
  pdftitle={Axion Strings AMR: Physics and Numerical Conventions},
  pdfauthor={Axion Strings AMR project}
}

\setlist[itemize]{leftmargin=1.4em}
\setlist[enumerate]{leftmargin=1.6em}

\newcommand{\psivec}{\psi}
\newcommand{\Rt}{\tilde{L}}

\title{\vspace{-1.5cm}\textbf{\Huge Axion Strings AMR}\\[0.3cm]\Large Physics and Numerical Conventions\\[0.15cm]\normalsize Network evolution only --- see the code for loop/collision test ICs}
\author{}
\date{\today}

\begin{document}
\maketitle
\vspace{-1.0cm}

\begin{note}
\textbf{Scope.} This is a physics reference, not a how-to: it specifies what equations are solved, how the
field is initialised and relaxed, how refinement is decided, and how every reported observable is defined
--- the detail needed to interpret a run's output correctly, not to launch one. For parameter names,
defaults and cluster mechanics, see the getting-started guide. Only the physical-string network pipeline is
covered (\texttt{ic\_mode = fourier\_relaxed} into the main evolution); the flat-space loop/collision test
ICs are a separate, simpler code path and are not described here.
\end{note}

\tableofcontents

\section{Field content, units, and normalisation}
\label{sec:units}

The field is a single complex scalar $\phi$ (the Peccei--Quinn field), written as two real components. Code
units fix
\[
v = 1, \qquad \lambda_0 = 1, \qquad \tau_0 = 1,
\]
where $v$ is the vacuum expectation value ($\langle|\phi|\rangle = v$), $\lambda_0$ the quartic coupling at
$\tau_0$, and $\tau$ is \emph{conformal} time throughout. The field is evolved in a rescaled, comoving form
\[
\psivec = R(\tau)\,\phi/v, \qquad \Pi \equiv \psivec' = \mathrm{d}\psivec/\mathrm{d}\tau,
\]
with $\psivec = \psivec_1 + i\psivec_2$; the four state variables are $(\psivec_1,\psivec_2,\Pi_1,\Pi_2)$.
Because $R(\tau)>0$ is real, $\arg(\psivec) = \arg(\phi)$ exactly --- the phase is unaffected by the
comoving rescaling, which is used throughout (string finding, masking) to work directly with the stored
state.

\textbf{Normalisation.} $\phi$ is canonically normalised as a \emph{complex} scalar: the kinetic and
gradient terms carry no extra factor of $\tfrac12$ (unlike a single real scalar). This was re-derived
directly from the equation of motion (Section~\ref{sec:eom}) rather than assumed, after finding an error of
this kind in an earlier document -- see \texttt{Energy.hpp}'s header for the check. The axion (phase) field
is $a = f_a\,\theta$, $\theta=\arg(\phi)$, with
\[
f_a = \sqrt{2}\,v = \sqrt{2} \quad \text{in code units.}
\]

\section{Background and equations of motion}
\label{sec:eom}

\subsection{Expansion law}
The scale factor is a pure power law in conformal time,
\[
R(\tau) = R_0\left(\frac{\tau}{\tau_0}\right)^{1/b_{\rm inv}}, \qquad b_{\rm inv} \equiv a_{\rm inv}-1,
\qquad R_0 = \frac{1}{b_{\rm inv}},
\]
so the single input \texttt{a\_inv} fixes the whole expansion history (radiation domination is
$a_{\rm inv}=2$). The physical Hubble rate is $H(\tau) = R'(\tau)/R(\tau)^2$, independent of everything
below --- only $\lambda(\tau)$, not $R(\tau)$, responds to the mass schedule.

\subsection{Mass schedule and the three regimes}
\label{sec:regimes}
\[
\lambda(\tau) = \lambda_0\left(\frac{R(\tau)}{R_0}\right)^{-2c}, \qquad m_r(\tau) = \sqrt{\lambda(\tau)}.
\]
A single exponent \texttt{c0} (with an optional single switch to \texttt{c1} at a chosen $\tau_{\rm sw}$,
$\lambda$ re-anchored there so it stays continuous, only its logarithmic slope changes) selects which of
three physically distinct schedules is being run:
\begin{itemize}
  \item \textbf{Physical strings}, $c=0$: $\lambda\equiv\lambda_0$ is exactly constant, so the string's
        comoving core width is fixed while the comoving horizon grows --- $m_r/H$ grows with $\tau$, which
        is the actual target dynamics (real global strings), and is why $\log(m_r/H)$ is the natural
        ``how far into the run'' coordinate used everywhere in this project.
  \item \textbf{Fat strings}, $c=1$: $\lambda\propto R^{-2}$, chosen so that $R(\tau)\,m_r(\tau)$ is exactly
        constant --- the comoving core width shrinks at exactly the rate the comoving horizon does, holding
        $m_r/H$ \emph{fixed}. Cores are artificially fat/cheap to resolve; used to relax a network onto its
        attractor fast (Section~\ref{sec:preevolution}) and, via a switch to $c=a_{\rm inv}$ mid-run, as a
        bridge into Moore's trick.
  \item \textbf{Moore's trick}, $c=a_{\rm inv}$: singular in the general box-planning formula below (the
        exponent $(a_{\rm inv}-1)/(a_{\rm inv}-c)$ diverges), handled by its own formulas
        (Section~\ref{sec:boxplan}). $H\!L$ (Hubble length over box size) falls as $1/\tau$ while resolution
        only improves, letting a single run reach a much larger $\log(m_r/H)$ than the physical schedule
        could afford at the same cost -- supporting evidence for the main physical-string result, not the
        headline measurement itself.
\end{itemize}

\subsection{Equation of motion}
The complex-scalar equation of motion actually integrated (one RHS, templated on nothing --- physical, fat
and Moore differ \emph{only} through $\lambda(\tau)$ above, never through separate code paths) is
\[
\Pi_i' = \nabla^2\psivec_i + c_{\rm curv}(\tau)\,\psivec_i - \frac{\lambda(\tau)}{2}\,\psivec_i\left(|\psivec|^2 -
R(\tau)^2\right), \qquad \psivec_i' = \Pi_i, \qquad i=1,2,
\]
with the Laplacian a comoving spatial derivative (fourth-order stencil) and
\[
c_{\rm curv}(\tau) = \frac{1-b_{\rm inv}}{b_{\rm inv}^2\,\tau^2}.
\]
$R(\tau)$, $\lambda(\tau)$ and $c_{\rm curv}(\tau)$ are evaluated on the host from $\tau$ and passed into
the kernel as plain numbers -- no parameter parsing or background evaluation inside the per-cell loop, so
the kernel is GPU-clean by construction.

\textbf{No Kreiss--Oliger dissipation.} $\sigma=0$ is the standing default: KO dissipation damps precisely
the high-$k$ modes the axion spectral index $q$ is measured from, so it is not added as a matter of course.
A nonzero $\sigma$, if ever required for stability, must be characterised (its effect on the spectrum
quantified) with the same care as the lattice spacing, not quietly switched on.

\section{Box planning}
\label{sec:boxplan}
Two physical targets, both evaluated at the run's final time $\tau_f$, fix the comoving box:
\begin{align*}
N_1 &= H(\tau_f)\,\Rt(\tau_f) &&\text{(Hubble patches per box length)}, \\
N_2 &= \frac{1}{\Delta(\tau_f)\,m_r(\tau_f)} &&\text{(grid points per core width)},
\end{align*}
with $N=N_1 N_2$ (comoving box length in finest-grid cells) the third input. For $c\neq a_{\rm inv}$,
\[
\tau_f = \left(\frac{N}{N_1 N_2}\right)^{\frac{a_{\rm inv}-1}{a_{\rm inv}-c}}, \qquad
\Rt = \frac{N_1}{R_0}\left(\frac{N}{N_1 N_2}\right)^{\frac{a_{\rm inv}-1}{a_{\rm inv}-c}}.
\]
For AMR, $N$/$N_2$ above refer to the \emph{finest} level; the base (level-0) grid is
$N/2^{\texttt{max\_level}}$. Each refinement level is only permitted once the current finest level's own
resolution has genuinely degraded to the $N_2$ target -- the level-addition schedule, spaced by
\[
\Delta\!\log(m_r/H)\ \text{per level} = \ln 2\cdot\frac{a_{\rm inv}-c}{1-c}
\]
(singular at $c=1$: a fat run's comoving core width never shrinks, so no further level is ever needed once
set up). Buffering between regrids (\texttt{amr.n\_error\_buf}) is sized so that the fastest possible
signal ($v=1$ in these units) cannot cross out of the tagged-plus-buffer region before the next regrid
check, given the fixed relation $\mathrm{d}t_{\rm level} = \texttt{dt\_multiplier}\times \mathrm{d}x_{\rm
level}$.

\section{Initial conditions and pre-evolution}
\label{sec:preevolution}

\subsection{Why relax at all}
Starting the physical-string ($c=0$) schedule directly from a raw random field would spend a large part of
the run just forming a realistic tangle, and the relaxation transient itself contaminates later
diagnostics. Instead, an over-seeded Gaussian random field is generated and evolved under a separate,
favourable schedule until it settles onto the network's own attractor, \emph{then} handed off into the main
($c=0$) run.

\subsection{Fourier-mode initial field}
$\psivec_1,\psivec_2$ are drawn as an independent Gaussian random field per Fourier mode (deterministic
per-mode hash of \texttt{(seed, component, i, j, k)}, not a stream-based RNG -- reproducible at fixed
\texttt{(seed, parameters, rank count)}), flat over $|k|\le k_{\max}$ and zero above, with
\[
k_{\max}\ (\text{in cells}) = \texttt{k\_max\_over\_mr}\times m_r(0)\times \mathrm{d}x,
\]
then rescaled in real space to hit the requested mean-square variance exactly. $k_{\max}$ is deliberately
\emph{not} restricted to the core scale: sub-core modes seed extra, doomed-to-annihilate vortex--antivortex
pairs (Kibble mechanism), and since annihilation only ever removes net winding, under-seeding can never
reach a dense target network -- $k_{\max}$ must be generous enough that the network is over-tangled at
$\tau_{\rm pre}=0$ and has genuine room to relax down.

\subsection{The pre-evolution background}
Relaxation runs on its own clock $\tau_{\rm pre}$ (starting at 0) and its own background, \emph{not} an
instance of Section~\ref{sec:eom}'s power law ($a_{\rm inv}=1$ is singular there). Switching to
$\tau_{\rm pre}$ turns a linear-in-cosmic-time $R$ into
\[
R_{\rm pre}(\tau_{\rm pre}) = e^{\tau_{\rm pre}}, \qquad
\lambda_{\rm pre}(\tau_{\rm pre}) = \frac{\gamma_{\rm pre}^2}{R_{\rm pre}(\tau_{\rm pre})^2}, \qquad
c_{\rm curv}^{\rm pre} = 1\ \text{(constant)},
\]
($R_0=\alpha=1$ WLOG, one physical input $\gamma_{\rm pre}$ remaining). This is exactly the fat-mode
$\lambda\propto R^{-2}$ shape (Section~\ref{sec:regimes}) under this new time gauge, so it reuses the same
RHS unchanged, and by construction holds
\[
\frac{m_r}{H}\bigg|_{\rm pre} = \gamma_{\rm pre} \quad \text{exactly constant throughout relaxation},
\]
letting the network converge onto a genuine attractor fast rather than chasing a moving target. Unless set
explicitly, $\gamma_{\rm pre} = 1/\mathrm{d}x_{\rm base}$ (resolves the core with $\sim\!1$ cell at
$\tau_{\rm pre}=0$) -- purely a starting-resolution choice; because $\lambda_{\rm pre}$ keeps decreasing
through relaxation, the comoving core width \emph{grows} over the course of relaxation, away from that
starting point.

\subsection{Stopping condition}
Relaxation stops the first time the measured $\xi$ (Section~\ref{sec:xi}, on the relaxation's own grid, at
an adaptive check cadence that tightens as $\xi$ approaches target) drops to \texttt{xi\_target}. $\xi$
generally decreases monotonically as the over-seeded network anneals.

\subsection{Handoff to the main run}
At the stopping instant $\tau_{\rm pre,end}$, the field is rescaled into the main run's own normalisation.
Writing $\phi=\psivec/R$ in either frame and requiring $\phi$ itself to be continuous across the switch
gives, with $\kappa \equiv R_{\rm main}(\tau_i)/R_{\rm pre}(\tau_{\rm pre,end})$,
\[
\psivec_i \to \kappa\,\psivec_i, \qquad
\Pi_i \to \kappa^2\,\Pi_i + \frac{\psivec_i}{R_{\rm pre}}\Big(R_{\rm main}'(\tau_i) - \kappa^2
R_{\rm pre}'(\tau_{\rm pre,end})\Big),
\]
an exact chain-rule identity (not an approximation): it matches $\psivec$ and $\mathrm{d}\psivec/\mathrm{d}\tau$
under the new background with no assumption beyond $\phi$'s own continuity. The main run's own clock then
begins at $\tau_i$ (\texttt{log\_mr\_over\_h\_i}), continuing from the same integrator state -- no restart
or checkpoint round-trip.

\subsection{Handoff smoothing}
\label{sec:handoffsmoothing}
The rescale above is exact for $\psivec,\Pi$, but $\lambda$ and $c_{\rm curv}$ themselves still jump
discontinuously at the handoff (the pre-evolution and main schedules are, in general, different curves in
$\lambda(\tau)$ that need not agree at whatever instant relaxation happens to stop). A discontinuous jump in
$\lambda$ is a discontinuous jump in the string core's natural width, and leaving the field's already-formed
spatial profile as-is while its equation of motion suddenly wants a different equilibrium width excites a
measurable transient: an oscillatory exchange between radial kinetic and potential energy, and a bump in
the axion spectrum near the old core scale that migrates towards the new one over several $\tau$. By
default (\texttt{handoff\_transition\_n\_periods = 3}), $\lambda$ and $c_{\rm curv}$ are instead blended
from their frozen pre-evolution values to the main schedule's own values via a smoothstep
($3t^2-2t^3$, value-continuous with zero slope at both ends, so the smoothing itself adds no new
discontinuity) over this many main-schedule core-oscillation periods, $2\pi/\sqrt{\lambda_{\rm
main}(\tau_i)}$. Set to $0$ to recover the instantaneous jump.

\section{Tagging and AMR refinement}
\label{sec:tagging}
The primary refinement criterion is exact and parameter-free: a cell is tagged if any of its three
low-index-corner plaquettes ($xy$, $yz$, $zx$) is pierced, i.e.\ the phase $\theta=\arg(\psivec)$ winds by a
full $\pm 2\pi$ around the plaquette's four corners (each consecutive step wrapped into $(-\pi,\pi]$ before
summing). This is the \emph{same} routine the $\xi$ diagnostic uses to count strings
(Section~\ref{sec:xi}) -- one implementation, not two copies that could silently drift apart. A level may
only be created once $\log(m_r/H)$ has actually crossed its own threshold (Section~\ref{sec:boxplan}); no
refinement happens during relaxation at all (relaxation's own clock is not $\log(m_r/H)$ in the main
schedule's sense, and holding $m_r/H$ fixed throughout relaxation means no new level is ever due there
regardless).

Two further, optional per-cell triggers exist -- a Laplacian threshold on $\psivec_1,\psivec_2$
(\texttt{gradient\_threshold}) and a radial-amplitude-gradient threshold
(\texttt{radial\_\allowbreak gradient\_\allowbreak threshold}) -- both \emph{off by default}: calibrating
either against an isolated static core was tried and found to flag essentially the whole domain once
applied to a real, dense network's field statistics, not just cores, so neither is trusted as a default.

\textbf{Refinement-systematics control.} \texttt{force\_full\_refinement} (off by default) bypasses both
the criteria above and the schedule-gating, tagging every cell at every level unconditionally -- not a
production setting, but a way to run an otherwise-identical pair of simulations, one adaptively refined and
one uniformly refined to the same finest resolution, to check directly whether refinement itself is biasing
any observable.

\section{Masking (screening)}
\label{sec:masking}
The axion phase rate $\theta' = (\psivec_1\Pi_2-\Pi_1\psivec_2)/|\psivec|^2$ genuinely diverges near a string
core ($|\psivec|\to0$), so any spectrum or energy built from it must exclude the core region before
transforming or summing. Two, largely independent mechanisms carry this out, one per quantity type,
each applied at exactly one place -- not inlined at several call sites, so an index or threshold error
cannot silently diverge between them.

\textbf{Energies} (Section~\ref{sec:energies}: $\rho_{\rm tot}$, axion and radial kinetic/gradient/mass) are
screened by a single multiplicative weight $w$, $\rho^{\rm scr}=w\,\rho^{\rm unscr}$, also used to correct
the effective point count in any spatial average over masked data ($n_{\rm unmasked}$ below):
\begin{itemize}
  \item \textbf{None}: $w=1$ identically -- the required unscreened pass every screened/unscreened pair
        needs (Sections~\ref{sec:energies},~\ref{sec:spectrum}) and the tension's core+tail term.
  \item \textbf{A} (smooth): $w = \big(1+r/f_a\big)^2$, $r\equiv|\phi|-v$ the radial deviation from vacuum
        (Section~\ref{sec:energies}'s own $r$), $f_a=\sqrt2\,v$ -- no hard cutoff, but a substantial,
        continuously-varying suppression near a core ($w\approx0.09$ exactly at $|\phi|=0$, rising to $1$ in
        the far field).
  \item \textbf{B} (top-hat): $w=1$ where $|\psivec|/R$ is at or above a runtime \texttt{threshold} (default
        $0.8$, not a settled choice -- $0.9$, $0.95$ have also been tried without reconciling on one value),
        $0$ below it.
\end{itemize}

\textbf{The axion spectrum} (Section~\ref{sec:spectrum}) instead masks by construction, filling the FFT
buffer with
\[
\dot a_{\rm scr} \propto
\begin{cases}
\theta' = (\psivec_1\Pi_2-\Pi_1\psivec_2)/|\psivec|^2 & \text{None (unscreened),} \\
\psivec_1\Pi_2-\Pi_1\psivec_2 & \text{A (bare numerator, no division),} \\
\theta'\ \text{with the scheme-B top-hat above} & \text{B.}
\end{cases}
\]
Scheme A's bare numerator is $=|\psivec|^2\theta'$ -- also a smooth, no-hard-cutoff suppression, but a
different one, specific to this construction, not literally $w\cdot\theta'$ with the same $w$ used for
energies above. \texttt{compute\_spectrum} requires scheme B specifically (checked at startup) -- the
persisted spectrum must be the genuinely screened field, so scheme A's spectrum branch, while implemented,
is never actually reachable in the saved \texttt{axion\_spectrum.dat} output.

\section{Observables}
\label{sec:observables}

\subsection{String density \texorpdfstring{$\xi$}{xi}}
\label{sec:xi}
Each pierced plaquette represents $\sim\!\tfrac23\,\mathrm{d}x$ of string length (a \emph{statistical}
correction for random orientation relative to the lattice -- exact for a network of many randomly-oriented
strings, not for a single deterministically axis-aligned one). Summing across AMR levels by comoving length
(never by raw cell/plaquette count, which would overweight finer levels),
\[
\ell_{\rm comoving} = \frac{2}{3}\sum_\ell N_{p,\ell}\,\mathrm{d}x_\ell, \qquad
\xi = \frac{\ell_{\rm comoving}}{\Rt^3}\cdot\frac{(a_{\rm inv}-1)^3}{a_{\rm inv}^2}\cdot\tau^2.
\]
$N_p$ (``plain'') counts pierced plaquettes; $N_{p,\rm weighted}$ sums $|{\rm winding}|$ over them -- most
strings are singly wound, so the two agree closely, and a divergence between them is itself a diagnostic.

\subsection{Energies}
\label{sec:energies}
The regular (non-singular) total physical energy density, from the canonically normalised $\phi$
Lagrangian directly (not the polar axion+radial decomposition below, which is individually singular at the
core):
\[
\rho_{\rm tot} = \frac{1}{R^4}\Big[\,|\Pi - (R'/R)\psivec|^2 + |\nabla\psivec|^2 +
\tfrac14\lambda\big(|\psivec|^2-R^2\big)^2\,\Big].
\]
The phase (axion) and amplitude (radial/Higgs) modes are also reported individually, each the genuine
energy carried by that one degree of freedom alone (not ``total minus the other''):
\[
\rho_a^{\rm kin} = \tfrac12\dot a^2,\qquad \rho_a^{\rm grad} = \tfrac12(\nabla a)^2, \qquad a = f_a\theta,
\]
with $\dot a = f_a\theta'/R$ (cosmic time from conformal) and $\nabla a$ picking up one power of $1/R$
(physical), not two (a pure phase has no amplitude to redshift). The radial energies use the exact
orthogonal decomposition of $|\dot\phi|^2$/$|\nabla\phi|^2$ into radial and tangential parts (a Pythagorean
identity, not an approximation): $r\equiv|\phi|-v$,
\[
\rho_r^{\rm kin} = \dot r^2, \qquad \rho_r^{\rm grad} = (\nabla r)^2, \qquad
\rho_r^{\rm mass} = \tfrac14\lambda\big(|\phi|^2-v^2\big)^2.
\]
Every energy is reported both screened and unscreened (Section~\ref{sec:masking}).

\subsection{String tension}
\label{sec:tension}
Two variants, both normalised by $\ell_{\rm comoving}$ above, from the composite (all-level) coverage sums
$n_{\rm total}$/$n_{\rm unmasked}$ (comoving \emph{volumes} once levels are combined, not raw counts):
\[
\mu_{\rm core} = \frac{R^2\big[n_{\rm total}\,\rho_{\rm tot}^{\rm unscr} - n_{\rm unmasked}\,\rho_{\rm
tot}^{\rm scr}\big]}{\ell_{\rm comoving}}, \qquad
\mu_{\rm core+tail} = \mu_{\rm core} + \frac{R^2\,n_{\rm unmasked}\big[\rho_{a,\rm scr}^{\rm grad} -
\rho_{a,\rm scr}^{\rm kin}\big]}{\ell_{\rm comoving}}.
\]
$\mu_{\rm core}$ is expected to converge to a resolution-independent \emph{constant}, but has no simple
closed form: it depends on both the near-core profile (not universal) and the masking threshold's exact
radius. $\mu_{\rm core+tail}$ is expected to \emph{diverge logarithmically} with the box/horizon size --
the long-range Goldstone tail is universal, and its coefficient is calculable directly from the axion
gradient energy density above: at large $\rho$ from an isolated winding-1 core, $|\nabla\theta|^2=1/\rho^2$,
so
\[
\rho_a^{\rm grad}(\rho) = \tfrac12 f_a^2/\rho^2 = 1/\rho^2 \ \ (v=1) \quad\Longrightarrow\quad
\mu_{\rm tail} = 2\pi\!\int_{1/m_r}^{L}\!\rho\,\frac{\mathrm{d}\rho}{\rho^2} = 2\pi\ln(m_r L).
\]
Taking $L\sim1/H$ (the natural IR cutoff for a single string in an expanding background) gives the
closed-form prediction
\[
\mu_{\rm core+tail} \approx 2\pi\,\ln(m_r/H)
\]
directly in terms of the reported $m_r/H$ column. This is the \emph{isolated single-string} estimate; a
real network's nearby strings partially screen each other's long-range tail, so a measured value somewhat
below this is expected, not a discrepancy.

\subsection{Spectrum}
\label{sec:spectrum}
The masked $a_{\rm dot}=f_a\theta'/R$ field (Section~\ref{sec:masking}) is Fourier-transformed
(\texttt{amrex::FFT::R2C}, level-0 only -- established policy, not yet extended across the hierarchy) and
binned into shells by integer mode number $p=\lfloor|\mathbf{p}|+0.5\rfloor$, dropping modes outside the
inscribed sphere ($p\ge N/2$). Reported alongside the bare shell average and mode index is the physical
$k/H$ for that mode,
\[
k/H = \frac{2\pi\,p}{\Rt\,H(\tau)},
\]
the $x$-axis a spectral index $q$ is actually fit against. Screened and unscreened spectra are both saved,
from two independent FFT passes (scheme B and None), so the impact of screening itself is directly
visible.

\subsection{String velocities}
\label{sec:velocity}
Evaluated at the corners of every pierced plaquette (global-string equilibrium profile coefficients
$c_1=0.41222$, $e_1=-0.025763$), R-independent by construction (a physical velocity cannot depend on the
comoving rescaling used to evolve the field -- verified as an exact algebraic cancellation, not assumed):
\[
\gamma^2v^2 = \frac{|\dot\phi|^2}{2m_r^2c_1^2}\left(1-\frac{e_1|\phi|^2}{c_1^3}\right) -
\frac{e_1\,[\mathrm{Re}(\phi^*\dot\phi)]^2}{4m_r^2c_1^5}.
\]

\end{document}
